%!TEX root = ../dokumentation.tex

\chapter{Annahmen und Risiken}

Gemäß Marchewka werden in diesem Abschnitt zentrale Annahmen dokumentiert,
auf denen die Planung des Projekts basiert, sowie wesentliche Risiken
identifiziert, die den Projekterfolg beeinflussen können. Risiken werden dabei
als unsichere Ereignisse oder Bedingungen verstanden, die Auswirkungen auf
Scope, Zeit, Kosten oder Qualität des Projekts haben können.

Die Eintrittswahrscheinlichkeit der Risiken wird qualitativ bewertet:

\begin{itemize}
    \item \textbf{Gering} – unwahrscheinlich, aber möglich
    \item \textbf{Mittel} – realistisches Risiko im Projektverlauf
    \item \textbf{Hoch} – wahrscheinlich ohne geeignete Gegenmaßnahmen
\end{itemize}


\section{Projektdefinition}

In der Phase der Projektdefinition werden Zielsetzung, Scope,
Rahmenbedingungen sowie die organisatorischen Grundlagen des
Projekts festgelegt. Diese Phase bildet die Basis für die weitere
Detailplanung und Umsetzung.

\subsection{Annahmen}

\begin{xltabular}{\textwidth}{@{}X X@{}}
    \hline
    \textbf{Annahme} & \textbf{Begründung} \\
    \hline

    Das Management unterstützt die Einführung einer zentralen
    Lernplattform strategisch und organisatorisch. &
    Voraussetzung für Ressourcenbereitstellung und
    projektweite Entscheidungsfähigkeit. \\

    Die strategischen Ziele der Organisation im Bereich
    Personalentwicklung und Weiterbildung bleiben während
    der Projektlaufzeit stabil. &
    Grundlage für den definierten MOV und die
    Projektzielsetzung. \\

    Die identifizierten Stakeholder sind bereit,
    aktiv an der Projektdefinition und den
    Planungsaktivitäten mitzuwirken. &
    Notwendig für Anforderungsdefinition und
    Akzeptanz des Projekts. \\

    Die organisatorischen Rahmenbedingungen
    (z.\,B. Budgetrahmen und Projektpriorität
    im Portfolio) bleiben während der
    Projektinitialisierung stabil. &
    Grundlage für die initiale Planung von
    Zeit, Ressourcen und Budget. \\

    Der kalkulierte administrative Aufwand im
    Schulungsmanagement von etwa 10 Mio. €
    pro Jahr bildet eine realistische Basis für
    die Wirtschaftlichkeitsbetrachtung. &
    Wird zur Bewertung des wirtschaftlichen
    Nutzens der Plattform im Business Case
    verwendet. \\

    \hline
\end{xltabular}


\subsection{Risiken}

\begin{xltabular}{\textwidth}{@{}X c X@{}}
    \hline
    \textbf{Risiko} & \textbf{Eintrittswahrscheinlichkeit} & \textbf{Strategie zur Minimierung oder Bewältigung} \\
    \hline

    Unklare oder widersprüchliche Erwartungen verschiedener Stakeholder führen zu einem unklaren Projektumfang und zusätzlichen Abstimmungsaufwänden in der Planungsphase. &
    Mittel &
    Durchführung einer strukturierten Stakeholderanalyse sowie moderierter Abstimmungsworkshops zur frühzeitigen Klärung der Erwartungen. \\

    Fehlende oder unzureichende Unterstützung durch das Top-Management kann zu eingeschränkter Ressourcenverfügbarkeit und einer niedrigen Priorität des Projekts im Portfolio führen. &
    Gering &
    Frühzeitige Einbindung des Sponsors sowie regelmäßige Statuskommunikation und Management-Reviews. \\

    Änderungen strategischer Prioritäten innerhalb der Organisation während der Projektinitialisierung können eine Anpassung oder Reduktion des Projektumfangs erforderlich machen. &
    Mittel &
    Klare Verankerung des Projekts im Projektportfolio sowie regelmäßige Abstimmung mit strategischen Entscheidungsgremien. \\

    Der organisatorische Veränderungsaufwand der Einführung einer zentralen Lernplattform wird unterschätzt, was zu geringer Akzeptanz und Verzögerungen bei der Umsetzung organisatorischer Prozesse führen kann. &
    Mittel &
    Frühe Einbindung von HR sowie strukturierte Change-Management-Aktivitäten und transparente Kommunikation des Nutzens. \\

    Unrealistische Erwartungen an die wirtschaftlichen Effekte des Projekts können zu einer negativen Bewertung des Projekterfolgs im Vergleich zum MOV führen. &
    Gering &
    Verwendung konservativer Annahmen im Business Case und transparente Definition der Erfolgsmessung. \\

    \hline
\end{xltabular}



\section{Detailplanung und Architekturdesign}

Diese Phase umfasst insbesondere die Analyse der Anforderungen,
die Definition der Zielarchitektur, das Infrastrukturdesign sowie
die Auswahl einer geeigneten Lernplattform.

\subsection{Annahmen}

\begin{xltabular}{\textwidth}{@{}X X@{}}
    \hline
    \textbf{Annahme} & \textbf{Begründung} \\
    \hline

    Die relevanten Fachbereiche stellen Anforderungen vollständig und
    fristgerecht bereit. &
    Voraussetzung für Stakeholder-Workshops und die Konsolidierung der Anforderungen. \\

    Eine geeignete Lernplattform kann im Markt identifiziert werden,
    die die funktionalen und technischen Anforderungen erfüllt. &
    Grundlage für Plattformauswahl und Architekturdesign. \\

    Die bestehende Systemlandschaft ermöglicht eine Integration
    über standardisierte Schnittstellen. &
    Voraussetzung für die Definition der Integrationsarchitektur. \\

    Unternehmensweite Architektur- und Sicherheitsrichtlinien sind bekannt
    und können im Architekturdesign berücksichtigt werden. &
    Notwendig für Architekturreview und Freigabe. \\

    Zentrale Stakeholder (HR, IT-Architektur, Betriebsrat) stehen während
    der Planungsphase für Entscheidungen zur Verfügung. &
    Erforderlich für Architekturreviews und Governance-Freigaben. \\

    Die organisatorischen Zielsetzungen des Projekts bleiben während der
    Planungsphase stabil. &
    Vermeidung grundlegender Scope-Änderungen. \\

    \hline
\end{xltabular}


\subsection{Risiken}

\begin{xltabular}{\textwidth}{@{}X c X@{}}
    \hline
    \textbf{Risiko} & \textbf{Eintrittswahrscheinlichkeit} & \textbf{Strategie zur Minimierung oder Bewältigung} \\
    \hline

    Unvollständige oder widersprüchliche Anforderungen aus verschiedenen Fachbereichen führen zu Verzögerungen im Architekturdesign und zusätzlichem Abstimmungsbedarf. &
    Hoch &
    Durchführung strukturierter Anforderungsworkshops sowie Priorisierung und Validierung der Anforderungen mit den Stakeholdern. \\

    Die geplante Zielarchitektur ist nur eingeschränkt mit der bestehenden IT-Systemlandschaft kompatibel, was zusätzlichen Integrationsaufwand verursacht. &
    Mittel &
    Frühe Analyse der bestehenden Systemlandschaft sowie Durchführung mehrerer Architekturreviews. \\

    Keine verfügbare Plattform erfüllt alle definierten Anforderungen vollständig, wodurch zusätzlicher Customizing- oder Integrationsaufwand entsteht. &
    Mittel &
    Durchführung eines Proof-of-Concept und strukturierte Bewertung der Plattformalternativen. \\

    Verzögerungen bei Architekturfreigaben durch Governance-Gremien oder Betriebsrat führen zu einer Verschiebung der Implementierungsphase. &
    Gering &
    Frühzeitige Einbindung der relevanten Entscheidungsträger sowie regelmäßige Abstimmungstermine. \\

    Der Integrationsaufwand zwischen Plattform und bestehenden Systemen wird unterschätzt, wodurch Termin- und Kostenüberschreitungen entstehen können. &
    Mittel &
    Durchführung einer detaillierten Schnittstellenanalyse sowie technischer Voruntersuchungen in der Planungsphase. \\

    \hline
\end{xltabular}



\section{Systemkonfiguration und Integration}

Diese Phase umfasst die Installation der Plattform, die technische
Grundkonfiguration sowie die Integration mit bestehenden HR-
und Unternehmenssystemen.

\subsection{Annahmen}

\begin{xltabular}{\textwidth}{@{}X X@{}}
    \hline
    \textbf{Annahme} & \textbf{Begründung} \\
    \hline

    Die ausgewählte Lernplattform kann innerhalb der vorgesehenen
    Infrastrukturumgebungen installiert werden. &
    Voraussetzung für die technische Implementierung. \\

    Die erforderlichen technischen Ressourcen und IT-Infrastruktur
    stehen während der Implementierungsphase zur Verfügung. &
    Notwendig für Installation, Konfiguration und Integration. \\

    Stammdaten aus HR-Systemen können für Benutzer- und
    Organisationsstrukturen genutzt werden. &
    Grundlage für Benutzer- und Rollenmodell der Plattform. \\

    Die benötigten Schnittstellen zu HR- und Reporting-Systemen
    sind technisch realisierbar. &
    Voraussetzung für Integration und automatisiertes Reporting. \\

    Die Datenmigration kann mit ausreichender Datenqualität
    durchgeführt werden. &
    Erforderlich für korrekte Benutzer- und Kursinformationen. \\

    Fachbereiche unterstützen die Definition von Rollen,
    Organisationsstruktur und Lernkategorien. &
    Notwendig für die Konfiguration der Plattformstruktur. \\

    \hline
\end{xltabular}


\subsection{Risiken}

\begin{xltabular}{\textwidth}{@{}X c X@{}}
    \hline
    \textbf{Risiko} & \textbf{Eintrittswahrscheinlichkeit} & \textbf{Strategie zur Minimierung oder Bewältigung} \\
    \hline

    Technische Integrationsprobleme mit bestehenden HR-Systemen führen zu Verzögerungen bei der Systemintegration und zusätzlichem Entwicklungsaufwand. &
    Mittel &
    Frühe Definition der Schnittstellenarchitektur sowie iterative Integrationstests während der Implementierung. \\

    Unzureichende Datenqualität bei der Migration bestehender Schulungs- und Benutzerinformationen führt zu fehlerhaften Datenstrukturen und eingeschränkter Systemnutzung. &
    Mittel &
    Durchführung von Datenbereinigung, Validierung und Testmigrationen vor der produktiven Datenübernahme. \\

    Fehlkonfiguration des Rollen- und Berechtigungsmodells führt zu Sicherheitsrisiken oder eingeschränkter Nutzung der Plattform. &
    Gering &
    Validierung der Berechtigungsstruktur gemeinsam mit Fachbereichen und IT-Security. \\

    Performanceprobleme bei hoher Nutzerzahl führen zu eingeschränkter Systemverfügbarkeit und geringer Nutzerzufriedenheit. &
    Gering &
    Einrichtung einer skalierbaren Infrastruktur sowie Durchführung von Last- und Performancetests. \\

    Verzögerungen bei der Bereitstellung der Infrastruktur durch interne IT-Abhängigkeiten verschieben die Installations- und Integrationsarbeiten. &
    Mittel &
    Frühzeitige Abstimmung mit IT-Operations und verbindliche Infrastrukturplanung. \\

    Geringe Akzeptanz der Plattformkonfiguration bei Fachbereichen führt zu Anpassungsbedarf nach der Implementierung. &
    Mittel &
    Iterative Abstimmung mit Fachbereichen und Durchführung von Pilotphasen. \\

    \hline
\end{xltabular}



\section{Rollout und Betrieb}

\subsection{Annahmen}

\begin{xltabular}{\textwidth}{@{}X X@{}}
    \hline
    \textbf{Annahme} & \textbf{Begründung} \\
    \hline
    -- & -- \\
    \hline
\end{xltabular}

\subsection{Risiken}

\begin{xltabular}{\textwidth}{@{}X c X@{}}
    \hline
    \textbf{Risiko} & \textbf{Eintrittswahrscheinlichkeit} & \textbf{Strategie zur Minimierung oder Bewältigung} \\
    \hline
    -- & -- & -- \\
    \hline
\end{xltabular}